% !TEX program = xelatex
\documentclass[12pt]{article}

% Pacchetti fondamentali
\usepackage{fontspec}
\usepackage{polyglossia}
% NOTA: NON caricare bidi esplicitamente, polyglossia lo fa da solo.

% Imposta la lingua principale
\setmainlanguage{italian}
% Definisci le altre lingue che userai
\setotherlanguages{arabic, hebrew, hindi, thai, japanese, chinese}

% Definizione dei font per ogni script
\newfontfamily\arabicfont[Script=Arabic, Scale=1.1]{Noto Sans Arabic}
\newfontfamily\hebrewfont[Script=Hebrew, Scale=1.1]{Noto Sans Hebrew}
\newfontfamily\devanagarifont[Script=Devanagari, Scale=1.1]{Noto Sans Devanagari}
\newfontfamily\thaifont[Script=Thai, Scale=1.1]{Noto Sans Thai}
\newfontfamily\japanesefont[Script=CJK, Language=Japanese]{Noto Sans CJK JP}
\newfontfamily\chinesefont[Script=CJK, Language=Chinese Simplified]{Noto Sans CJK SC}

\begin{document}

\section*{Test Multilingua}

\subsection*{Italiano (LTR)}
Questo è un testo in italiano. La lingua principale.

\subsection*{Arabo (RTL)}
\textarabic{مرحبا بالعالم! هذا اختبار للنص العربي.}

\subsection*{Ebraico (RTL)}
\texthebrew{שלום עולם! זהו מבחן לטקסט בעברית.}

\subsection*{Hindi (LTR)}
\texthindi{नमस्ते दुनिया! यह हिंदी पाठ का एक परीक्षण है।}

\subsection*{Thai (LTR)}
\textthai{สวัสดีชาวโลก! นี่คือการทดสอบข้อความภาษาไทย}

\subsection*{Giapponese (LTR)}
\textjapanese{こんにちは世界！これは日本語のテストです。}

\subsection*{Cinese Semplificato (LTR)}
\textchinese{你好世界！这是一个中文测试。}

\end{document}
